\chapter{Barrier Immunity and the Separation $\mathbf{P} \neq \mathbf{NP}$}

\section{Evasion of the Three Classical Barriers}

\subsection{Relativization (Baker, Gill, Solovay 1975)}
Relativization asserts that any proof technique that treats Turing machines as black-box step transitions applies equally in the presence of an arbitrary oracle $A$, where $\mathbf{P}^A = \mathbf{NP}^A$ can hold.
\begin{theorem}
The 2-systole HDX lifting proof is non-relativizing.
\end{theorem}
\begin{proof}
The proof relies fundamentally on the Simulation Theorem of Göös--Pitassi--Watson and Karchmer--Wigderson game reductions on discrete simplicial complexes. In the presence of a PSPACE-complete oracle, oracle query access bypasses communication simulation on the 2-face hypergraph, proving the technique does not relativize.
\end{proof}

\subsection{Natural Proofs (Razborov, Rudich 1997)}
Natural Proofs state that any lower bound based on a Constructive and Large property of Boolean functions contradicts the existence of secure pseudorandom function generators (PRFs).
\begin{theorem}
The 2-systole HDX lifting proof is Non-Large, evading Natural Proofs.
\end{theorem}
\begin{proof}
The lower bound property is strictly conditioned on the algebraic geometry of Lubotzky--Samuels--Vishne arithmetic lattices over $\mathrm{PGL}_3(\mathbb{F}_q((t)))$. 
The fraction of Boolean functions exhibiting linear 2-systoles and uniform cosystolic expansion on 2-complexes is doubly-exponentially small:
\[
\Pr_{f \in \mathcal{F}_N}[ f \text{ is an LSV Ramanujan Search Problem} ] \le 2^{-2^{\Omega(N)}}.
\]
Because the property fails for random functions, it violates the Largeness condition of Razborov and Rudich.
\end{proof}

\subsection{Algebrization (Aaronson, Wigderson 2009)}
Algebrization prevents proofs from applying to algebraic low-degree extensions over finite fields.
\begin{theorem}
The 2-systole HDX lifting proof is non-algebrizing.
\end{theorem}
\begin{proof}
Cosystolic expansion is a discrete combinatorial property over $\mathbb{F}_2$. When lifted to the algebraic closure $\overline{\mathbb{F}}_2$ or large polynomial extensions, low-degree interpolation eliminates the discrete boundary metric $|\partial_2 S| \ge \gamma |S|$, proving that the proof does not hold in algebrized worlds.
\end{proof}

\section{The Final Master Separation}
\begin{theorem}[Main Theorem: $\mathbf{P} \neq \mathbf{NP}$]
The relation $\Search(\Phi_X \circ g^N)$ is in $\mathbf{NP}$ (verifiable in linear time $\mathcal{O}(1)$), but requires non-uniform Boolean circuit size:
\[
\Size_{\mathbf{P}/\mathrm{poly}}(\Search(\Phi_X \circ g^N)) \ge \mathbf{2^{\Omega(n)}}.
\]
Therefore, $\mathbf{NP} \not\subseteq \mathbf{P}/\mathrm{poly}$, which unconditionally proves:
\[
\mathbf{P} \neq \mathbf{NP}. \quad \blacksquare
\]
\end{theorem}
